% Historical deployment protocol retained for provenance.
% Current complete-cohort selection is documented in benchmark/README.md.
\paragraph{Scope of this protocol record}
The deployment details below describe the historical evaluation archive, including the partial ZMFF and ReDiffuse cohorts. The current complete comparison excludes both until all four datasets have been completed and validated. IFCNN is evaluated by the recorded sequential pairwise adaptation. Its generic implementation can also process multiple inputs jointly, but that joint variant is not evaluated by the current records.

\section{Baseline deployment protocol}
\label{supp:baselines}

The seven newly deployed comparison methods (all evaluated in the main paper; ZMFF and ReDiffuse on Mobile Depth and Middlebury only) use the code revisions listed in Table~\ref{tab:deploy}. Pretrained methods use the released checkpoints in inference mode without retraining or fine-tuning; ZMFF uses its prescribed per-pair zero-shot optimization without pretrained weights. Pairwise methods fuse each stack iteratively, $F_1=g(X_0,X_1)$ and $F_k=g(F_{k-1},X_k)$, where the images are ordered by a natural sort of their file names, so that numeric focal-plane names are compared by value. For GIFNet, DeFusion, CSNet and ZMFF the intermediate result is stored as an 8-bit image between steps. GIFNet, ReDiffuse and ZMFF fuse the luminance channel and treat chroma with a separate rule defined by their official code, while DeFusion, CSNet, MFIF-GAN and VAEEDOF process RGB directly.

The six classical spatial-domain methods (GFF, IFM, QB, BFMF, GFDF and MGFF) were run from the MATLAB code released by their authors, without modification, under GNU Octave~10.3 with the image package. GFF, IFM and QB accept the whole stack in one call; the other three fuse pairs iteratively in numeric file order, with the fused image quantized to 8 bit between pairs, as in their published usage. Octave and MATLAB give pixel-identical outputs for these codes on the scenes we compared (one scene per benchmark, all six methods). As a check of the deployment itself, we also fused the five Lytro pairs of MFIFB\footnote{\url{https://github.com/xingchenzhang/MFIFB}} and compared the results with the images that MFIFB releases for the same pairs: all six methods agree with the published images to 40~dB or better (Table~\ref{tab:supp_mfifb}), the residual difference being the JPEG compression of the published files. MGFF was designed for infrared/visible pairs, and on stacks it degrades as it is applied iteratively; Table~\ref{tab:supp_mfifb} validates the single-pair implementation; the stack results expose the degradation under iterative application (Section~4.2 of the main paper).

VAEEDOF is trained for stacks of seven images. Its official \texttt{predict()} call returned broken output for longer stacks: on Middlebury Adirondack with 15 images the output mean was 48.0, against 119.0 for the input images, and these outputs were discarded. For stacks with more than seven images we therefore fuse the first seven images, then fuse the previous result with the next six images, and pad the last group by repeating its final image. This is a sequential chain of the seven-image model and not a single native pass over all images. On Adirondack the chained output has mean 119.4 and standard deviation 47.6, close to the 119.5 and 46.6 obtained from the first seven images alone, but this check does not establish that the chain preserves accuracy on other scenes. FILM was excluded. Its inference path requires a 768-dimensional text feature for each image pair, computed from ChatGPT descriptions and BLIP-2 features. The repository does not ship these features, none of our four datasets provides them, and the restormer block hard-codes $288\times384$ feature maps, so the stacks cannot be processed at their native resolution without resizing. ReDiffuse was run from a runtime copy because the release omits \texttt{B\_Conv.py}, so its compiled bytecode was loaded directly. This was validated by a strict weight load and by reproducing the released Lytro sample (PSNR 71.7 dB against the released output). The luminance output of ReDiffuse on Mobile Depth \emph{balls} is darker than the input image (mean 94.6 against 130.1). A likely cause is the per-image min-max stretch of the official \texttt{tensor2img} function. Because a full run of ReDiffuse over the four benchmarks would take several GPU-days, it is evaluated on Mobile Depth and Middlebury only. For ZMFF, the official code is not bit-reproducible on our GPU: two official runs with the same seed differ by a mean absolute luminance difference of 4.1 on a 0 to 255 scale, and our port differs from the official output by 4.5, which we treat as the same level of run-to-run variation.

The order of the images matters for every method that fuses a stack pair by pair. On FlyingThings3D the file names are focal distances (\texttt{10.0}, \texttt{16.43}, \dots, \texttt{100.0}) and on Road-MF they are integers without leading zeros, so sorting them as strings, as the MATLAB \texttt{dir()} function and several original scripts do, interleaves near and far focal planes. We therefore regenerated the results of the pairwise and transform-domain baselines with the images sorted by the numeric value of their names, using the original code with no other change. The transform-domain methods were rerun with the original MATLAB code in GNU Octave on FlyingThings3D and Road-MF; on the validation scene (Mobile Depth \emph{balls}, where string and numeric order coincide) the Octave results of CVT, DWT, NSCT and DCT are pixel-identical to the stored MATLAB results, so the stored results are kept for Mobile Depth and Middlebury. DTCWT differs in Octave (mean absolute difference 2.8 gray levels) and is regenerated in MATLAB. The learned pairwise methods (IFCNN, SwinFusion, MUFusion, SwinMFF, DDBFusion, CCSR-Net, MCCSR-Net) were rerun with their original scripts on all four benchmarks. On Mobile Depth, where the order is unchanged, the reruns differ from the stored results by a mean absolute difference of 1 to 4 gray levels, which we attribute to differences in software versions (a run with reversed image order differs by 7 to 32 gray levels); to keep each method consistent, all four benchmarks of these methods are taken from the reruns. MFF-GAN, U2Fusion and SDNet are TensorFlow~1 models; we redeployed them from their official repositories (commits \texttt{2f674c8}, \texttt{0a58821} and \texttt{bf6a2ee}) with TensorFlow~2.13 running the original graphs through its \texttt{tf.compat.v1} interface (TensorFlow~1.15 does not support our GPUs), verified each on the official two-image demonstration, and applied them in the same image-by-image manner. Their luminance on Mobile Depth differs from the stored benchmark results by a mean of 1.7 (SDNet), 9.1 (U2Fusion) and 19.0 (MFF-GAN) gray levels; for MFF-GAN the stored results were produced with a locally modified test script, which we could not reproduce, so we report the results of the unmodified official code. For three datasets the stored DWT results of the benchmark were byte-identical to the NSCT results (126 files). Re-running the respective original transform scripts reproduced the stored NSCT outputs and the Road-MF DWT outputs exactly. We therefore retained NSCT and regenerated DWT using its original implementation. StackMFF was rerun with its officially released checkpoint. The results of StackMFF V2, Helicon Focus 8 and Zerene Stacker are taken from the StackMFF benchmark. For the transform-domain methods, the exact result roots and regenerated outputs are recorded with the evaluation configuration; all metrics are recomputed from the result images. 

Helicon Focus 8 outputs are slightly smaller than the ground truth (for example 956 by 538 against 960 by 540). We align each output to the ground truth from its original size, using ORB features (2000 keypoints) and a RANSAC homography with a reprojection threshold of 5 px. The image is interpolated once, and uncovered borders are filled by edge replication. When the homography has fewer than 10 inliers or an implausible scale, the translation estimated by phase correlation is used instead. For Helicon Focus 8 the median homography scale is about 1.000 (cropping changes output dimensions without requiring proportional rescaling) and the per-dataset median inlier counts lie between 1248 and 1581. All other methods have median phase-correlation shifts of at most 0.22 px per benchmark and are evaluated without alignment.
